\section{Conclusion}
\label{sec:conclusion}

This paper presented \textbf{MPTA-Repair}, an automated workflow for
multi-property and multi-counterexample clock-bound repair of
industrial timed automata. To overcome the limitations of isolated
single-trace fixes, the method accumulates diagnostic traces and
exploits explicit relations among properties, counterexamples, and
repairable bounds to guide a MaxSMT-based search. Each synthesized
candidate is validated through full-property regression and an
admissibility check based on untimed functional equivalence
\cite{kolbl2019clock}. This design ensures that accepted repairs are
system-level timing corrections that preserve the intended operational
behavior.

Evaluations across benchmark and industrial datasets show that
MPTA-Repair outperforms the iterative baseline, particularly in
industrial systems with dense property interactions. The failure
analysis further indicates that unsuccessful repairs often stem from
the intrinsic limits of a strictly parametric clock-bound repair space.

Future work will primarily focus on extending the automated repair space
beyond numerical clock bounds to include broader structural
modifications, such as altering synchronization labels and redesigning
discrete control logic. We also plan to investigate the theoretical and
algorithmic relationship between candidate repair generation and
admissibility verification. By formally integrating behavioral
admissibility constraints into the early stages of search-space
formulation, we aim to proactively prune functionally invalid candidates
and thereby improve computational efficiency for large-scale industrial
applications.
